\section{A root-centered equation in thirteen positive integers}
\label{sec:positive13}

The construction in this section has thirteen strictly positive integer
unknowns, exact total degree $114$, and $83$ live binary arithmetic
operations: $46$ multiplications, $19$ additions, and $18$ subtractions.
Squaring and multiplication by a fixed scalar are charged. Intermediate
registers are polynomial abbreviations and do not add unknowns.

The new step merges the first Pell root and the auxiliary Pell ordinate
using a centered two-root norm. Keeping the ordinary strong coordinate
allows direct recovery of both deleted roots. This gives a degree
improvement over the pinned thirteen-positive constructions of degree
$126$ in $91$ operations and degree $132$ in $83$ operations. These are
comparisons with the named repository constructions, without a claim of
global minimality or literature priority.

\subsection{The fixed equation}

Use the entire compact44 native tuple, with
$\beta=2^{590069}-1$. The four payloads and their exact hashes are fixed by
\texttt{COEFFICIENT\_BINDING.json}; the aliases are
$\beta\mapsto\beta$, $K_0\mapsto K$, $M_C\mapsto M_C$, and
$M_F\mapsto M_{F,\mathrm{source}}$.
The last coefficient is already shifted:
$M_{F,\mathrm{source}}=M_{F,\mathrm{native}}+2^{590069}-1$.
The whole tuple has $295388$ raw bytes. Its included proof supplies the
same native blank-machine compiler interface, with the ordinary strong
witness obtained from a canonical normalized one by
$i_{\mathrm{ordinary}}=\Delta i_{\mathrm{normalized}}$.

Supply precisely
\[
 J,F,\alpha,z,f,i,T,s,w,\eta,\zeta,\gamma,t
 \quad\in\mathbb Z_{>0}.
\]
Define the following polynomial abbreviations:
\begin{align*}
 q&=\beta J+1,& C&=q-F-\alpha,\\
 X&=wq,&Y&=sq^3,&E&=XY,& k&=\eta+\zeta,\\
 c&=kY+\eta,&a&=E+Y,&H&=4a+3,&\Delta&=a^2+H,\\
 \ell&=EkY,& A_1&=1+\ell(\ell+k),&G&=1+\Delta c^2,&b_0&=X+ac.
\end{align*}
The packing and transport expressions are
\begin{align*}
 R&=\bigl((q-1)(q-F)+\alpha\bigr)(q^2-1)
       +(M_C+qM_F)J,\\
 N_t&=(K_0+w)C+q-F-z(q-1).
\end{align*}
For the absorbed main root put
\begin{align*}
 U&=4q^4,& r&=k-R-1,\\
 D_*&=Ub_0+H(U\gamma-r),&
 N_{\mathrm{main}}&=D_*^2-U^2G.
\end{align*}
The strong and auxiliary quantities are
\begin{align*}
 I&=ic^2,&Q&=\Delta(f^2-1),&d&=Q-1,&N_s&=I^2-d,\\
 V&=c(Tf-1)-Rf^2,&E_{\mathrm{aux}}&=QV^2-1,\\
 x&=t+\ell,&\delta&=x^2-A_1,&
 N&=(d\delta-E_{\mathrm{aux}})^2-4dA_1E_{\mathrm{aux}}.
\end{align*}
The proposed equation is
\begin{equation}
 \boxed{\displaystyle
 P_{13}=N_t(1+N_{\mathrm{main}})N_s(1+N)-1=0.}
 \label{eq:p13-main}
\end{equation}
Every displayed power has a fixed nonnegative integer exponent.
The included division-free JSON circuit expands every abbreviation.
The coordinate $z$ here is the old positive \emph{product} transport
coordinate. It equals $z_{\mathrm{shift}}+1$ for the coordinate used in a
shifted equation $(K_0+w)C=F+z_{\mathrm{shift}}(q-1)$.

With the fixed tuple, ordinary mathematics formalizable in ZFC proves
\[
 \exists\boldsymbol v\in\mathbb Z_{>0}^{13}\;
 P_{13}(\boldsymbol v)=0
 \quad\Longleftrightarrow\quad
 \text{there is a finite syntactic ZFC refutation}.
\]
Negating this equivalence gives the required equivalence between the
absence of positive zeros and $\operatorname{Con}(\mathrm{ZFC})$.
The native compiler and the literal K215 calibration are inherited
theorems; no new end-to-end proof-assistant certificate is claimed.

\subsection{The centered two-root identity}

For arbitrary indeterminates $d,x,A,e$ there are exact polynomial identities
\begin{align}
 (d(x^2-A)-e)^2-4dAe
 &=\bigl(d(x^2+A)-e\bigr)^2-4d^2x^2A \notag\\
 &=d^2(x^2-A)^2-2d(x^2+A)e+e^2.
 \label{eq:p13-norm-identity}
\end{align}
The first expression saves one multiplication in the evaluator. If
$d,x,A$ are positive integers and it vanishes, the middle expression gives
the rational number
\[
 \tau_0=\frac{d(x^2+A)-e}{2dx},\qquad \tau_0^2=A.
\]
A rational square root of an integer is an integer. Therefore
$\tau_0\in\mathbb Z$ and
\begin{equation}
 d(x-\tau_0)^2=e.
 \label{eq:p13-second-root}
\end{equation}
This direct argument allows coincident square classes. It requires no
unproved independence of formal radical signs and no auxiliary rational
Pell saturation lemma.

The particular shift $x=t+\ell$ gives
\begin{equation}
 \delta=t^2+\ell(2t-k)-1.
 \label{eq:p13-centering}
\end{equation}
This cancellation lowers its degree from the apparent $22$ to $12$.
The converse will prove positivity of the new witness
$t=\tau-\ell+y$; positivity is not assumed to be preserved by an arbitrary
affine change of an old tuple.

\subsection{The unit guard and the two restored roots}

At a positive zero of \eqref{eq:p13-main}, each of the four integer factors
is $1$ or $-1$. Since $4\mid U$, the residue of $N_{\mathrm{main}}$ modulo
$4$ is the square $D_*^2$. The residue of $N$ modulo $4$ is also a square.
Neither residual can be $-2$. Thus
$N_{\mathrm{main}}=N=0$ and their two factors are $+1$.

For every integer $a$, $\Delta=a^2+4a+3$ is $0$ or $3$ modulo $4$.
Hence $Q=\Delta(f^2-1)$ is $0$ or $1$ modulo $4$.
It follows that $N_s=1+I^2-Q$ is $0$, $1$, or $2$ modulo $4$ and cannot
be $-1$. Consequently
\begin{equation}
 Q=I^2,\qquad N_t=1,\qquad N_{\mathrm{main}}=N=0.
 \label{eq:p13-guard-conclusions}
\end{equation}
Here $c>1$ and $I=ic^2>1$, so $Q>1$ and $d>0$. In particular the
apparently troublesome case $f=1$ is excluded by the equations themselves.
Also $x=t+\ell>0$ and $A_1>0$.

If $V=0$, then $E_{\mathrm{aux}}=-1$, and
\[
 N=(d\delta+1)^2+4dA_1>0,
\]
a contradiction. Therefore $V\ne0$ and $E_{\mathrm{aux}}>0$.
Applying \eqref{eq:p13-norm-identity} and
\eqref{eq:p13-second-root} gives integers
$\tau_0$ and $y_0=x-\tau_0$ with
\[
 \tau_0^2=A_1,\qquad dy_0^2=E_{\mathrm{aux}}.
\]
Thus the positive integers $\tau=|\tau_0|$ and $y=|y_0|$ satisfy
\begin{equation}
 \tau^2-\ell(\ell+k)=1,\qquad
 QV^2-(Q-1)y^2=1,\qquad
 (ic^2)^2=\Delta(f^2-1).
 \label{eq:p13-restored-roots}
\end{equation}
The identity
$\ell(\ell+k)=(XY^2)(XY^2+1)k^2$ identifies the first original Pell norm.
Together with \eqref{eq:p13-guard-conclusions}, these are precisely the
ordinary-strong fourteen-coordinate absorber equations, with all retained
coordinates still positive. Neither an index quotient nor its sign has
yet been assumed.

\subsection{Early index recovery and restoration of the compiler}

The original absorber proof is supplied verbatim as
\texttt{upstream/FOURTEEN\_ABSORPTION\_PROOF.md}; its full source identity
is recorded in the manifest. We restate its proof order because the early
index bound is necessary for the ordinary strong-rank lemma.

The equality $N_t=1$ implies $C>0$: if $C\le0$, positivity of $K_0+w$ gives
\[
 N_t\le q-F-z(q-1)=1-F-(z-1)(q-1)\le0.
\]
The authentic mask bounds then give
\begin{equation}
 q\ge16,\quad 3q<R<q^4-2,\quad E\ge q^4>R+2,\quad
 kY<c<k(Y+1),\quad a+2<2E.
 \label{eq:p13-outer-bounds}
\end{equation}
From $N_{\mathrm{main}}=0$, $g=D_*/U$ is a rational square root of the
integer $G$, hence an integer. It is positive. Indeed, $H<5a$, $b_0>akY$,
$r<k$, and $\gamma>0$ imply
\[
 D_*>Ub_0-Hk>ak(UY-5)>0.
\]
Thus $g^2-\Delta c^2=1$. Reduction of $D_*=Ug$ modulo $U$ gives
\begin{equation}
 U\mid H(k-R-1).
 \label{eq:p13-absorbed-congruence}
\end{equation}

Put $A_0=a+2$ and $P_0=2XY^2+1$. The usual positive Pell sequences give
indices $p,n\ge1$ such that
\[
 g=\chi_{A_0}(p),\quad c=\psi_{A_0}(p),\qquad
 \tau=\chi_{P_0}(n),\quad k=2\psi_{P_0}(n).
\]
Since $P_0>A_0$ and $c>k$, monotonicity implies $p>n$.
The first recurrence gives $k\equiv2n\pmod{4XY^2}$, and
$U\mid4XY^2=4ws^2q^7$. Moreover $H$ is odd and $H\equiv3\pmod q$;
if $3\mid q$, then $H\equiv3\pmod{27}$. Hence
$\gcd(U,H)\le3$. Defining $\mu=U/\gcd(U,H)$, we obtain
\[
 \mu\ge\frac{4q^4}{3}>R+1,\qquad
 2n\equiv R+1\pmod\mu,
 \qquad n\ge\frac{R+1}{2}\ge25.
\]
The lower bound precedes recovery of $p=R$. It gives $p\ge26$, and Pell
growth therefore proves
\begin{equation}
 c>A_0\Delta^2,\qquad c>2p,\qquad c>2R.
 \label{eq:p13-rank-hypotheses}
\end{equation}
For example,
$c\ge(2A_0-1)^{p-1}\ge A_0^{25}>A_0(A_0^2-1)^2$;
also $c\ge\psi_{A_0}(2)=2A_0>2R$.

The ordinary strong-rank lemma in the supplied
\texttt{ORDINARY\_BLANK\_PROOF.md}, Section~4, now applies. It yields
\begin{equation}
 \begin{gathered}
 f=\chi_{A_0}(m),\quad p\mid m,\quad c\mid m,\quad m\ge c>2p,\\
 I=\Delta\psi_{A_0}(m),\qquad c^2\mid(f^2-1).
 \end{gathered}
 \label{eq:p13-strong-rank}
\end{equation}
This lemma treats the possible nonfundamentality of $A_0$ in the underlying
Pell group. It is used with the verified strict hypothesis
\eqref{eq:p13-rank-hypotheses}; it is not replaced by informal cancellation
of $\Delta$ in the strong equation.

The auxiliary equation in \eqref{eq:p13-restored-roots}, the congruence
$V\equiv-c\pmod f$, and the odd abscissa quotient identity give an odd
auxiliary index $j_0$ with $j_0\equiv\pm p\pmod m$.
On the other hand $Q\equiv0\pmod c$, $f^2\equiv1\pmod c$, and
$V\equiv-R\pmod c$ give $j_0\equiv\pm R\pmod c$.
The signed index theorem from the ordinary source, Section~5, applies
because $0<2p<m$. Since $c\mid m$ and $0<p,R<c/2$, it follows that
$p=R$. This step permits either sign of $V$.

Now $n<p=R$, so $|2n-R-1|<R<\mu$, and the same congruence forces
$2n=R+1$. Consequently $r=k-2n$ is divisible by $4XY^2$, hence by both
$E$ and $U$. Set
\[
 h=\frac rE,\qquad \gamma_{\mathrm{old}}=\gamma-\frac rU.
\]
The first is a positive integer because $\psi_{P_0}(n)>n$ for $n\ge25$.
The second equals $(g-b_0)/H$. Since $n\ge2$, $k\ge4P_0>X$ and $c>X$;
furthermore $g>(a+1)c$, so $g-b_0>c-X>0$.
Thus every original ordinary compiler coordinate has been restored with
its required positive sign. The inherited compiler soundness gives the
blank K215 halt and hence a finite ZFC refutation.

\subsection{Constructive completeness}

Conversely, start with a canonical positive ordinary compiler witness of
a blank K215 halt. Canonical normalized witnesses supply these by
$i_{\mathrm{ordinary}}=\Delta i_{\mathrm{normalized}}$.
Its index theorem gives $p=R=2n-1$, so $r=k-2n$ is a positive multiple
of $4XY^2$ and therefore of $U$. Define
\[
 \gamma=\gamma_{\mathrm{old}}+\frac rU,
 \qquad t=\tau-\ell+y.
\]
Both are positive integers. In particular
$\tau^2=1+\ell(\ell+k)>\ell^2$ implies $\tau>\ell$.
Retain all other supplied coordinates. Then $x=\tau+y$,
$A_1=\tau^2$, and $E_{\mathrm{aux}}=dy^2$, so
\eqref{eq:p13-norm-identity} gives $N=0$.
Also $D_*=Ug$, giving $N_{\mathrm{main}}=0$, while $N_t=N_s=1$
are unchanged. Therefore \eqref{eq:p13-main} holds.
This proves both existential directions, with no additional cofinality
condition on an auxiliary index.

\subsection{Exact degree and finite verification}

All witnesses have degree one and the four ports degree zero. Put
\[
 q_0=\beta J,\qquad E_0=wsq_0^4,\qquad
 c_0=(\eta+\zeta)sq_0^3,\qquad \ell_0=E_0c_0.
\]
Their degrees are $1,6,5,11$, respectively. In the inner norm,
\[
 \deg d=14,\quad \deg V=7,\quad\deg E_{\mathrm{aux}}=28,
 \quad\deg A_1=22,\quad\deg\delta=12.
\]
The three terms of the final expression in
\eqref{eq:p13-norm-identity} have degrees $52,64,56$.
Only the middle one reaches degree $64$, with homogeneous part
\begin{equation}
 N_{64}=-4E_0^6c_0^4T^2f^6.
 \label{eq:p13-inner-head}
\end{equation}
In the efficient evaluator the squared core has degree at most $56$;
its correction $-4dA_1E_{\mathrm{aux}}$ has exact degree $64$.

The main residual has the exact identity
\begin{align*}
 N_{\mathrm{main}}
 &=U^2\bigl[2ac(X+H\gamma)+(X+H\gamma)^2-Hc^2-1\bigr]\\
 &\quad-2UHr(X+ac+H\gamma)+H^2r^2.
\end{align*}
Its three displayed terms have degrees $26$, at most $25$, and at most
$20$, so
\[
 (N_{\mathrm{main}})_{26}=128q_0^8E_0^2c_0\gamma.
\]
The other factors have leaders
\[
 (N_s)_{22}=i^2c_0^4,\qquad
 (N_t)_2=w(q_0-F-\alpha)-zq_0.
\]
Thus $\deg P_{13}\le2+26+22+64=114$.
On the full diagonal, where every witness is replaced by one
indeterminate, the four leading coefficients are
\[
 -2,\qquad 256\beta^{19},\qquad16\beta^{12},\qquad-64\beta^{36}.
\]
Their product is $2^{19}\beta^{67}\ne0$ on the fixed coefficient slice.
This proves exact degree $114$ without treating the ports as additional
variables or using a numerical experiment as an upper-degree proof.

The primary circuit uses $83$ binary gates, all live. The independent
standard-library checker imports no builder. It authenticates the whole
tuple and all source pins; independently verifies $41$ named formulas
against the literal circuit; proves the norm, centering, and main
cancellation identities; checks the universal weighted degree bound;
and evaluates six complete polynomial rays over three finite fields
using the actual fixed tuple. Normal and optimized Python runs agree.
These finite checks verify transcription, fixed data, and degree. They do
not replace the unbounded mathematical compiler and calibration proofs.

The original source50 absorber, ordinary strong-rank proof, compact native
adapter, and source59 comparison are preserved in the included upstream
directory. Their exact source paths, commits, Git blob identities and
SHA256 digests are recorded in \texttt{SOURCE\_PINS.json}. The complete
formal-assistant dependency tree of the inherited K215 endpoint is not
claimed to have been rebuilt in this work.
